\documentclass[10pt,a4paper]{amsart}
\usepackage[margin=19mm]{geometry}
\usepackage{amsmath,amssymb,mathtools}
\usepackage[scr=rsfs,cal=euler]{mathalfa}
\usepackage{xfrac}
\usepackage[unicode,hidelinks,bookmarks=false]{hyperref}
\newcommand{\ie}{i.e.\ }
\newcommand{\cf}{cf.\ }
\newcommand{\resp}{respectively\ }
\newcommand{\Subsection}{\subsection}
\providecommand{\nobreakdash}{\nobreak\hbox{-}}
\setlength{\emergencystretch}{2em}
\multlinegap0pt
\usepackage[shortlabels]{enumitem}
\usepackage{dsfont}
\usepackage{amssymb}
\usepackage{tikz}
\usepackage{float}
\newtheorem{lem}{Lemma}
\newtheorem{pro}{Proposition}
\newtheorem{cor}{Corollary}
\title{\bf Co-edge-regular graphs with four eigenvalues and unbounded coherent rank}
\author{Edwin R. van Dam, Hong-Jun Ge, Jack H. Koolen}
\date{Source: arXiv:2606.19981v1 (CC BY 4.0)}
\begin{document}
\maketitle

\noindent\emph{Source and licence.} \bf Co-edge-regular graphs with four eigenvalues and unbounded coherent rank. Version arXiv:2606.19981v1 (CC BY 4.0), distributed by arXiv under the Creative Commons Attribution 4.0 International licence, http://creativecommons.org/licenses/by/4.0/, as recorded in the arXiv licence metadata for this exact version and shown on the abstract page https://arxiv.org/abs/2606.19981v1. Full licence notice: \texttt{licenses/arxiv-2606.19981-CC-BY-4.0.txt}.\par\smallskip

\noindent\textit{Editorial scope.} Counted unit: the article's cor cor:tpl-level (source0.tex lines 1010--1020). The article's complete original proof of it is delivered with it (source0.tex lines 1022--1117, one \texttt{proof} environment of 96 lines). No mathematics is written, repaired or re-derived: the statement and the proof are the article's own lines, in the article's own notation, and no retained line is substituted or re-flowed.  Retained uncounted as the source-local context the proof needs: lines 602--629; lines 825--833; lines 912--935.  Nothing was omitted from any retained span, so the package carries no omission record.  No retained line contains a citation, so the wrapper carries no bibliography and no bibliography entry.  No retained line contains a figure environment, an image-inclusion command or drawing code, so no image is delivered and the package declares no asset dependency.  Editorial formatting only, in addition: an AMS \texttt{amsart} preamble that restates the article's own declarations and macros used by the retained lines, namely the environments \texttt{lem}, \texttt{pro}, \texttt{cor} on separate counters, together with the standard packages those lines need.  The retained environments print in this document as Lemma 1, Pro 1, Corollary 1 in order of appearance, whereas the article prints the same items with the numbering of its own full document.  The complete native source of the article as distributed in the arXiv e-print for this exact version is bundled unchanged as \texttt{sources/203183/source0.tex}.
\begin{lem}\label{lem:basic-structure}
For $0\le j\le i\le \lfloor q/2\rfloor$, the following hold in $\Gamma_{i,j}$.
\begin{enumerate}[label=\textup{(\roman*)}]
\item The induced subgraph on $P_\Pi^+\cup P_\Pi^-$ is complete.
\item For every line label $L\in\mathcal L$,
\[
L^-\setminus P_\Pi^-=\{x^-:x\in L\setminus P_\Pi\},
\qquad
|L^-\cap P_\Pi^-|=|L\cap P_\Pi|.
\]
Moreover, for any two line labels $L_1,L_2\in\mathcal L$,
\[
|L_1^-\cap L_2^-|=|L_1\cap L_2|,
\qquad
|L_1^-\cap L_2^-\cap P_\Pi^-|=|L_1\cap L_2\cap P_\Pi|.
\]
\item If $x\in P\setminus P_\Pi$, then
\[
N_{\Gamma_{i,j}}(x^+)\cup\{x^+\}
=
N_{\Gamma_{i,j}}(x^-)\cup\{x^-\}.
\]
In particular,
\[
\bigl|N_{\Gamma_{i,j}}(x^+)\cup\{x^+\}\bigr|=2(q+1)\beta+2.
\]
\end{enumerate}
\end{lem}

\begin{pro}\label{prop:same-layer}
For any distinct \(x,y\in P\), the following hold in \(\Gamma_{i,j}\).
\begin{enumerate}[label=(\roman*)]
\item If $x\nsim y$ in $G$, then $x^+\nsim y^+$ and
$\mu_{\Gamma_{i,j}}(x^+,y^+)=2\mu_G(x,y).$
\item If $x\sim y$ in $G$, then $x^+\sim y^+$ and
$\lambda_{\Gamma_{i,j}}(x^+,y^+)=2\lambda_G(x,y)+2.$
\end{enumerate}
\end{pro}

\begin{pro}\label{prop:special-values}
Let $1\le t\le i$.
The following estimates hold in $\Gamma_{i,i}$.
\begin{enumerate}[label=\textup{(\roman*)}]
\item
$0\le
\lambda_{\Gamma_{i,i}}(u_t^+,u_t^-)
-
\bigl(2|P_\Pi|-2+2(q+1-t)(\beta-\alpha)\bigr)
\le 2t^2.$
\item
$0\le
\lambda_{\Gamma_{i,i}}(u_t^+,u_{\widetilde t}^-)
-
\bigl(2|P_\Pi|-2+2(t+1)(\beta-\alpha)\bigr)
\le 2(q-t)^2.$
\item If $w\in \ell_{1,1}\setminus\bigl(\{u_1\}\cup \ell_{q,1}\bigr),$
then
$
0\le
\lambda_{\Gamma_{i,i}}(u_1^+,w^-)-(2|P_\Pi|-2)
\le 2(q+1)^2.$
\end{enumerate}
\end{pro}

\begin{cor}\label{cor:tpl-level}
Assume that
$\beta\ge |P_\Pi|+(q+1)^2.$
If $\Lambda(G)=\gamma$, then
\[
\Lambda(\Gamma_{i,i})\ge \gamma+2i+2
\qquad\text{for } i<q/2.
\]
Moreover,
$\Lambda(\Gamma_{\lfloor q/2\rfloor,\lfloor q/2\rfloor})\ge \gamma+q+1.$
\end{cor}

\begin{proof}
Let
$\Lambda_0=
\{2\lambda_G(x,y)+2:x\sim y\text{ in }G\}.$
By Proposition~\ref{prop:same-layer}(ii), every value in $\Lambda_0$ occurs as a $\lambda$-value in $\Gamma_{i,i}$, and $|\Lambda_0|=\gamma$.
If $x\sim y$ in $G$, then $x$ and $y$ lie on a unique line of $\mathcal S$.
This line contributes $\beta-1$ common neighbours.
Each of the remaining $q$ lines through $x$ meets each of the remaining $q$ lines through $y$ in at most one point, so
\[
\beta-1\le \lambda_G(x,y)\le \beta-1+q^2.
\]
Therefore
\begin{equation}\label{eq:lambda0-interval}
\Lambda_0\subseteq [2\beta,2\beta+2q^2].
\end{equation}

Choose $w$ as in Proposition~\ref{prop:special-values}(iii).
Then
\[
\lambda_{\Gamma_{i,i}}(u_1^+,w^-)
\le 2|P_\Pi|-2+2(q+1)^2
<2\beta,
\]
so this value lies strictly below all values in $\Lambda_0$.

Put
$D=\beta-\alpha.$
Since $\beta\ge |P_\Pi|+(q+1)^2$ and $\alpha\le q$, we have
\begin{equation}\label{eq:D-large}
D>(q-1)^2.
\end{equation}
For $1\le t\le i$, set
\[
A_t=\lambda_{\Gamma_{i,i}}(u_t^+,u_{\widetilde t}^-),
\qquad
B_t=\lambda_{\Gamma_{i,i}}(u_t^+,u_t^-).
\]
By Proposition~\ref{prop:special-values},
\[
2|P_\Pi|-2+2(t+1)D
\le A_t
\le
2|P_\Pi|-2+2(t+1)D+2(q-t)^2,
\]
and
\[
2|P_\Pi|-2+2(q+1-t)D
\le B_t
\le
2|P_\Pi|-2+2(q+1-t)D+2t^2.
\]
Using \eqref{eq:D-large}, these estimates imply, for $i<q/2$,
\[
A_1<A_2<\cdots<A_i<B_i<B_{i-1}<\cdots<B_1.
\]
Indeed, consecutive separations follow from the inequalities
\[
D>(q-t)^2,
\qquad
D>t^2,
\qquad
D\ge \frac{(q-i)^2}{q-2i},
\]
which are all consequences of \eqref{eq:D-large} in the indicated ranges.
Moreover,
$A_1>2\beta+2q^2,$
again by the lower bound for $A_1$ and the hypothesis on $\beta$.
Thus all $A_t$ and $B_t$ are strictly above the interval \eqref{eq:lambda0-interval}.
Finally, for any $y\in P\setminus P_\Pi$, Lemma~\ref{lem:basic-structure}~(iii) gives
$$\lambda_{\Gamma_{i,i}}(y^+,y^-)=2(q+1)\beta,$$
and this is strictly larger than $B_1$.
Hence, for $i<q/2$, the graph $\Gamma_{i,i}$ has at least
\[
\gamma+1+i+i+1=\gamma+2i+2
\]
distinct $\lambda$-values.

Now assume that $q$ is even and $i=q/2$.
The same estimates give
\[
A_1<\cdots<A_{i-1}<B_i<B_{i-1}<\cdots<B_1,
\]
with all these values above $\Lambda_0$, while $\lambda_{\Gamma_{i,i}}(u_1^+,w^-)$ is below $\Lambda_0$ and $\lambda_{\Gamma_{i,i}}(y^+,y^-)$ is above $B_1$.
This gives at least
\[
\gamma+1+(i-1)+i+1=\gamma+q+1
\]
distinct $\lambda$-values.
For odd \(q\), we have \(\lfloor q/2\rfloor<q/2\), so the first part gives
\[
\Lambda(\Gamma_{\lfloor q/2\rfloor,\lfloor q/2\rfloor})
\ge \gamma+2\lfloor q/2\rfloor+2
=\gamma+q+1.
\]
The even case was just proved, so the final assertion follows.
\end{proof}

\end{document}
